\medskip

On considère le cube ABCDEFGH qui est représenté ci-dessous.

\psset{unit=1cm,arrowsize=2pt 3}
\begin{pspicture}(9.3,9)
\pspolygon(0.5,1.3)(5.5,0.6)(5.5,6.3)(0.5,7)
\uput[dl](0.5,1.3){A} \uput[d](5.5,0.6){B} \uput[u](5.5,6.3){F} \uput[l](0.5,7){E} 
\psline(5.5,0.6)(8.2,1.7)(8.2,7.4)(5.5,6.3)
\uput[r](8.2,1.7){C} \uput[ur](8.2,7.4){G} 
\psline(8.2,7.4)(3.7,8.1)(0.5,7)
\uput[u](3.7,8.1){H}
\psline[linestyle=dashed](0.5,1.3)(3.7,2.4)(3.7,8.1)
\psline[linestyle=dashed](3.7,2.4)(8.2,1.7)
\uput[d](3.7,2.4){D}
\end{pspicture}


Dans le repère orthonormé $\left(\text{A}~;~\vect{\text{AB}}~;~\vect{\text{AD}}~;~\vect{\text{AE}}\right)$, on considère les points M, N et P de coordonnées :
\begin{center}M$\left(1~;~1~;~\dfrac34\right)$, \qquad N$\left(0~;~\dfrac12~;~1\right)$, \qquad P$\left(1~;~0~;~-\dfrac54\right)$\end{center}

Dans cet exercice, on se propose de calculer le volume du tétraèdre FMNP.

\medskip

\begin{enumerate}
\item Donner les coordonnées des vecteurs $\vect{\text{MN}}$ et $\vect{\text{MP}}$.
\item Placer les points M, N et P sur le cube ci-dessus.
\item Justifier que les points M, N et P ne sont pas alignés. 

Dès lors les trois points définissent le plan (MNP).
\item 
	\begin{enumerate}
		\item Calculer le produit scalaire $\vect{\text{MN}} \cdot \vect{\text{MP}}$, puis en déduire la nature du triangle MNP.
		\item Calculer l'aire du triangle MNP.
	\end{enumerate}
\item 
	\begin{enumerate}
		\item Montrer que le vecteur $\vect{n}(5~;~-8~;~4)$ est un vecteur normal au plan (MNP).
		\item En déduire qu'une équation cartésienne du plan (MNP) est $5x - 8y + 4z = 0$.
	\end{enumerate}	
\item On rappelle que le point F a pour coordonnées F(1~;~0~;~1). 

Déterminer une représentation paramétrique de la droite $d$ orthogonale au plan (MNP) et passant par le point F.
\item On note L le projeté orthogonal du point F sur le plan (MNP). 

Montrer que les coordonnées du point L sont: L$\left(\dfrac47~;~\dfrac{24}{35}~;~\dfrac{23}{35}\right)$.
\item Montrer que FL $= \dfrac{3\sqrt{105}}{35}$ puis calculer le volume du tétraèdre FMNP. 

On rappelle que le volume V d'un tétraèdre est donné par la formule :

\begin{center}$V = \dfrac13 \times  \text{aire d'une base} \times \text{hauteur associée à cette base}$.\end{center}

\end{enumerate}

\medskip
